% ============================================================================
%  Chapter 10. Conclusion
%
%  \input file. Requires natbib in the parent preamble.
%  Bibliography: accord/thesis/refs.bib
%
%  Short by design. Every number restated here is derived in Chapters 3 to 9.
%  No new claim, no new number.
% ============================================================================

\chapter{Conclusion}
\label{ch:conclusion}

\section{What Was Asked}
\label{sec:conc-asked}

A patient on eleven medications is discharged. A nursing note records that the
loop diuretic was held because the patient had become hypovolemic. The discharge
summary instructs that it be continued at the admission dose. The two documents
disagree, only one of them governs what happens next, and it is not the one
containing the reason the instruction is wrong.

This thesis asked whether that disagreement can be detected automatically. It
posed four questions: whether conditioning on note-type authority beats
symmetric inference, whether temporal metadata helps, whether published
guidelines can rank confirmed contradictions by danger, and whether weak
supervision from a language model can replace expert annotation at the scale a
detector requires. It built the pipeline the field would naturally build to
answer them: rule-based candidate mining, language model weak supervision, and a
purpose-built neural detector over a domain-pretrained clinical encoder.

\section{What Was Found}
\label{sec:conc-found}

The clinical problem is real. Four transparent rules over one institution's notes
surface $1{,}978$ candidates across $1{,}203$ admissions from a heart-failure
cohort of $9{,}783$, and a board-certified pharmacist confirms $12$ genuine
contradictions in a $299$-pair stratified sample, four of them rated $4$ or $5$
for danger. That is systematic, not anecdotal.

The pipeline fails at every stage.

Candidate mining enriches the pool to a $4.0\%$ base rate, which is a real gain
over an unfiltered background rate orders of magnitude lower, but it attaches
the wrong evidence to $41.4\%$ of action-conflict candidates, because the trigger
quantified over sets of medication mentions while the exemplar selector indexed
positionally, and because no field existed upstream in which the distinction
could be expressed. The annotator, the weak labeller, and the model all reasoned
about that evidence.

Weak supervision fails on calibration rather than on comprehension. The labeller
assigned the positive class to $84.9\%$ of $1{,}679$ candidates against a $4.0\%$
measured base rate, so at least $95\%$ of its positive labels are wrong. The
student then reproduced the teacher's per-rule predicted-positive rate to within
$1.5$ percentage points on three of four rules, while the true prevalences on
those rules differ by more than a factor of four, and on the one rule where the
teacher was genuinely selective the student regressed to its own prior instead
of inheriting the distinction.

Three architectural contributions do not survive an audit of the code that
produced their numbers. The authority parameter reported as discovered at
$(0.5, 1.0)$ was initialised at exactly $[0.5, 1.0]$, moved by $1.8\%$ and
$0.6\%$, sits on a branch a subsequent normalisation largely erases, and is
outperformed by the symmetric variant on both reported metrics. The auxiliary
rule head's $F_1$ of $0.9603$ is a copy of an input feature and collapses to
$0.0654$ and $0.0595$ on two transfer cohorts. The severity component has no
guideline corpus, receives no gradient, and emits scores spanning $0.398$ to
$0.610$, the range of an untrained sigmoid.

The detector has no rank signal. On the $299$ annotated pairs its AUROC is
$0.460$, with average precision at the no-skill line, Mann-Whitney $p = 0.64$,
and a bootstrap $95\%$ interval of $[0.287, 0.633]$; its fifteen highest-scored
pairs contain zero true contradictions; and at every threshold from $0.05$ to
$0.15$ its predicted-positive rate is exactly $1.000$. Restricted to the strata
that cannot exhibit the exemplar defect, AUROC is $0.5078$ with Mann-Whitney
$p = 0.953$, on six positives. Its Matthews correlation is $-0.007$; the
labeller it was distilled from scores $-0.044$. A general-purpose model prompted
zero-shot on the same evidence, with no training and no clinical components, is
the only system in the comparison above chance.

Of the four questions: RQ4 is answered negatively with a mechanism; RQ2 is
unaddressable in MIMIC-III, because \texttt{CHARTTIME} is null for every
discharge summary; and RQ1 and RQ3 were never tested, because in each case the
component that would have answered them could not have produced an informative
answer under any data.

\section{What It Means for Anyone Attempting This}
\label{sec:conc-means}

Three statements are worth carrying out of this document.

\paragraph{The binding constraint is upstream of the model.} A candidate pool at
$4.0\%$ precision determines what every downstream stage can achieve, and no
detector in this thesis adds lift over the rule that generated its candidates:
precision tracks each rule's base rate to within a tenth of a percentage point,
and on the one rule where the model exercises selectivity it selects wrongly.
Effort spent on architecture was effort spent on the wrong stage. Anyone
building this system should spend it on extraction, on evidence attribution, and
on candidate generation, and should validate each against human annotation
before a model exists.

\paragraph{The task is not the problem.} A zero-shot general-purpose model, on
the same $299$ pairs, with the same evidence, reaches $8.45\%$ precision at
$50\%$ recall with a positive Matthews correlation. That is not a usable
clinical system, and it should not be presented as one: eleven false alarms per
true contradiction, half the contradictions missed. But it is above chance,
which means the evidence contains extractable signal and the failure documented
here belongs to this pipeline rather than to the problem. Every hypothesis this
thesis set out to test remains open.

\paragraph{The intended intervention point does not exist.} The system is
triggered by the discharge summary, and the discharge summary is a retrospective
document that frequently is not finalised until the patient has gone. No alert
built on this trigger can prevent anything. Either detection moves to the
discharge order, which changes the task from cross-note textual inference toward
reconciliation against a structured order set, or the contribution is
retrospective audit and quality assurance and should be described that way. That
question should be settled before any of the modelling work is repeated.

\section{The Transferable Lesson}
\label{sec:conc-lesson}

The most durable result in this thesis is methodological, and it is not about
contradictions.

A system can produce a plausible headline number while every component that
supposedly produces it is either untrained, leaking, or reading the wrong
evidence. This one did. It reported a contradiction $F_1$, an auxiliary-task
$F_1$ of $0.9603$, an authority hierarchy described as discovered, a severity
score, an ablation table with error bars, and a $\kappa$ of $1.000$ offered as
evidence of annotation reliability. Behind those numbers: a parameter that never
moved from its initialisation, an auxiliary label written into the model's own
input, two output heads absent from the loss, a standard deviation over two
seeds plus a pasted literal compared against a reference row that the ablation
script does not run, and an agreement statistic between two quantities standing
in an identity relation. Every one of these passed internal review. None of them
survives its own first diagnostic.

The diagnostics are not sophisticated. Print a module manifest at training time
with parameter counts and loss-graph membership. Log
$\lVert\theta_{\text{final}} - \theta_{\text{init}}\rVert /
\lVert\theta_{\text{init}}\rVert$ for any parameter reported as learned. Assert
that an auxiliary head's target is not derivable from its inputs, and check by
transfer rather than by held-out accuracy. Report a majority-class baseline
beside every $F_1$. Print the number of distinct values a confidence field takes.
Print the empirical range of every score the model emits. Estimate the base rate
on a human-labelled sample \emph{before} generating weak labels, and print the
labeller's marginal against it. Assert group disjointness of splits, and assert
that the evaluation set is absent from training, by identifier, at load time.

None of these is novel. All of them are omitted routinely, including here. What
this thesis contributes is the demonstration that their absence is not a hygiene
issue but a load-bearing one: it is what allowed four unsupported claims to
reach a submission, and it is the reason a system with no measurable
discriminative ability could be described, in good faith, as working.

The corollary is uncomfortable and is the reason this document exists in its
present form. The failures documented here were not found by peer review, by a
held-out set, or by a novel evaluation. They were found by an adversarial audit
of the project's own code against the project's own claims, conducted after the
results were in. A pipeline that runs end to end, produces numbers, and is never
compared against its own specification will report the specification. Only
deliberate self-falsification surfaces the difference, and the discipline that
makes it routine is the thing this field is missing more than it is missing
another architecture.

\section{Closing}
\label{sec:conc-closing}

Dangerous medication contradictions between nursing notes and discharge
summaries are real, recurring, and worth detecting. The pipeline the field would
naturally build for them fails at every stage. The failures are independent of
one another, each has an identifiable mechanism, each has a detector costing a
few lines of code, and none of them was found until the project set out to
falsify itself.

That is the contribution. It is a negative result, it is measured four
independent ways, and it is offered so that the next attempt starts further
along than this one did.
